%!TEX program = xelatex
\documentclass[8.5pt,a4paper]{article}

\usepackage{fontspec}
\usepackage{polyglossia}
\setmainlanguage{russian}
\setotherlanguage{english}

\setmainfont{SF Pro Text}[
  BoldFont       = {SF Pro Text Semibold},
  ItalicFont     = {SF Pro Text Regular Italic},
  BoldItalicFont = {SF Pro Text Semibold Italic},
]
\setmonofont[
  Scale=0.80,
  Path=/Users/ivanm3/Library/Fonts/,
  UprightFont = JetBrainsMonoNerdFontMono-Regular.ttf,
  BoldFont    = JetBrainsMonoNerdFontMono-Bold.ttf,
]{JetBrainsMono NFM}

\usepackage[
  top=6mm, bottom=5mm,
  left=9mm, right=9mm,
  columnsep=5mm,
  textheight=284mm
]{geometry}
\usepackage{multicol}
\usepackage{xcolor}
\usepackage{enumitem}
\usepackage{listings}
\usepackage{microtype}
\usepackage{titlesec}
\usepackage{pagecolor}

% ── Palette ──────────────────────────────────────────────────────────────────
\definecolor{bg}{HTML}{0F1117}
\definecolor{border}{HTML}{30363D}
\definecolor{blue}{HTML}{58A6FF}
\definecolor{green}{HTML}{3FB950}
\definecolor{yellow}{HTML}{D29922}
\definecolor{red}{HTML}{F85149}
\definecolor{muted}{HTML}{8B949E}
\definecolor{fg}{HTML}{E1E4E8}
\definecolor{codefg}{HTML}{C9D1D9}
\definecolor{codebg}{HTML}{0D1117}

\pagecolor{bg}
\color{fg}

% ── Section headings ─────────────────────────────────────────────────────────
\titleformat{\section}[block]
  {\normalfont\fontsize{7.5}{9}\bfseries\color{blue}}
  {}{0em}{}[\vspace{-1pt}\textcolor{border}{\rule{\linewidth}{0.4pt}}]
\titlespacing{\section}{0pt}{4pt}{1.5pt}

% ── Code listings ─────────────────────────────────────────────────────────────
\lstdefinestyle{gocode}{
  language=Go,
  basicstyle=\ttfamily\fontsize{6.0}{7.5}\color{codefg},
  backgroundcolor=\color{codebg},
  frame=single, framesep=2pt, framerule=0.4pt,
  rulecolor=\color{border},
  breaklines=true, keepspaces=true, columns=flexible,
  xleftmargin=3pt, xrightmargin=3pt,
  aboveskip=1.5pt, belowskip=1.5pt,
  keywordstyle=\color{blue},
  stringstyle=\color{yellow},
  commentstyle=\color{muted},
  morekeywords={chan,go,select,defer,func,var,type,struct,interface,make,close,range,return,case,default,for,if,else,package,import,sync},
  literate={<-}{{\texttt{<-}}}2 {<}{{\texttt{<}}}1 {>}{{\texttt{>}}}1,
}

\lstdefinestyle{sqlcode}{
  language=SQL,
  basicstyle=\ttfamily\fontsize{6.0}{7.5}\color{codefg},
  backgroundcolor=\color{codebg},
  frame=single, framesep=2pt, framerule=0.4pt,
  rulecolor=\color{border},
  breaklines=true, keepspaces=true, columns=flexible,
  xleftmargin=3pt, xrightmargin=3pt,
  aboveskip=1.5pt, belowskip=1.5pt,
  keywordstyle=\color{blue}\bfseries,
  stringstyle=\color{yellow},
  commentstyle=\color{muted},
}

\lstset{style=gocode}

% ── Compact bullet list ───────────────────────────────────────────────────────
\newenvironment{clist}{%
  \begin{itemize}[leftmargin=9pt,itemsep=0pt,topsep=0.5pt,parsep=0pt,
    label=\textcolor{muted}{\fontsize{5}{5}\textbullet}]
    \setlength{\itemsep}{0pt}
}{\end{itemize}}

% ── Helpers ───────────────────────────────────────────────────────────────────
\newcommand{\hl}[2]{\textbf{\textcolor{#1}{#2}}}
\newcommand{\note}[1]{{\fontsize{6.5}{8}\color{muted} #1}}

\setlength{\parindent}{0pt}
\setlength{\parskip}{0.5pt}

% ─────────────────────────────────────────────────────────────────────────────
\begin{document}
\pagestyle{empty}
\enlargethispage{12mm}

% ══ TITLE ════════════════════════════════════════════════════════════════════
{\centering
  {\fontsize{12.5}{15}\bfseries\color{blue} Go Concurrency \& SQL — Cheat Sheet}\\[-1pt]
  {\fontsize{7}{8.5}\color{muted} Паттерны каналов · sync-примитивы · SQL задачи}\\[1pt]
  \textcolor{border}{\rule{\linewidth}{0.5pt}}\\[1.5pt]
}

\begin{multicols}{3}
\setlength{\columnsep}{5mm}

%──────────────────────────────────────────────────────
\section{Каналы: основы}

{\fontsize{7}{8.5} Небуферизованный — синхронный (блок до получателя). Буферизованный — асинхронный до заполнения.}
\begin{lstlisting}
ch := make(chan int)     // sync
ch := make(chan int, 10) // buffered
\end{lstlisting}

\hl{green}{Аксиомы:}
\begin{clist}
  \item send в \texttt{nil}/closed → \hl{red}{panic}
  \item recv из closed → zero\,+\,\texttt{false}
  \item recv/send в nil → блок навсегда
  \item \texttt{select}: nil-ветка никогда не выбирается
\end{clist}

%──────────────────────────────────────────────────────
\section{Pipeline}

{\fontsize{7}{8.5} Цепочка горутин: каждая читает \texttt{in}, пишет в \texttt{out}, закрывает после себя.}
\begin{lstlisting}
func stage(in <-chan int) <-chan int {
  out := make(chan int)
  go func() {
    defer close(out)
    for v := range in { out <- v*v }
  }()
  return out
}
// usage: src -> stage -> stage -> sink
\end{lstlisting}

%──────────────────────────────────────────────────────
\section{Fan-out / Fan-in}

\hl{yellow}{Fan-out}: N воркеров читают один канал.\\
\hl{green}{Fan-in}: слить N каналов в один.
\begin{lstlisting}
func merge(cs ...<-chan int) <-chan int {
  var wg sync.WaitGroup
  out := make(chan int)
  wg.Add(len(cs))
  for _, c := range cs {
    go func(ch <-chan int) {
      defer wg.Done()
      for v := range ch { out <- v }
    }(c)
  }
  go func() { wg.Wait(); close(out) }()
  return out
}
\end{lstlisting}

%──────────────────────────────────────────────────────
\section{Done-channel / Отмена}

\begin{lstlisting}
done := make(chan struct{})
defer close(done) // broadcast-сигнал

go func() {
  select {
  case out <- process():
  case <-done: return // leak-safe
  }
}()
\end{lstlisting}

\note{Предпочти \texttt{context.Context} в production-коде.}

%──────────────────────────────────────────────────────
\section{Timeout / Select}

\begin{lstlisting}
select {
case v := <-ch:
  handle(v)
case <-time.After(2 * time.Second):
  // timeout
case <-ctx.Done():
  return ctx.Err()
}
\end{lstlisting}

%──────────────────────────────────────────────────────
\section{Семафор (ограничение пула)}

\begin{lstlisting}
sem := make(chan struct{}, N)
for _, job := range jobs {
  sem <- struct{}{} // acquire slot
  go func(j Job) {
    defer func() { <-sem }()
    process(j)
  }(job)
}
\end{lstlisting}

%──────────────────────────────────────────────────────
\section{Mutex / RWMutex}

\begin{lstlisting}
var mu sync.RWMutex
// чтение (параллельно):
mu.RLock(); v := m[k]; mu.RUnlock()
// запись (эксклюзивно):
mu.Lock(); defer mu.Unlock()
m[k] = v
\end{lstlisting}

\begin{clist}
  \item \texttt{defer mu.Unlock()} сразу после \texttt{Lock()}
  \item Pointer receiver — мьютекс \textbf{нельзя копировать}
  \item \hl{red}{Не реентерабелен} — повторный Lock = дедлок
\end{clist}

%──────────────────────────────────────────────────────
\section{WaitGroup / Once / Atomic}

\begin{lstlisting}
var wg sync.WaitGroup
wg.Add(1)        // ДО go func
go func() {
  defer wg.Done(); work()
}()
wg.Wait()

var once sync.Once
once.Do(initDB)  // ровно 1 раз

var cnt int64
atomic.AddInt64(&cnt, 1)
v := atomic.LoadInt64(&cnt)
\end{lstlisting}

%──────────────────────────────────────────────────────
\section{Типичные ошибки}

\begin{clist}
  \item \hl{red}{Goroutine leak}: горутина заблокирована навсегда → done-ch или ctx
  \item \hl{red}{Data race}: r/w без sync → \texttt{go test -race}
  \item \hl{yellow}{wg.Add внутри go func} → Wait может выйти раньше
  \item \hl{yellow}{Value-receiver + Mutex} → копирует мьютекс
  \item \hl{yellow}{Закрывает получатель} → panic у второго отправителя
\end{clist}

%──────────────────────────────────────────────────────
\section{SQL: схема}

{\fontsize{6.5}{8}\ttfamily\color{codefg}
employees(id, name, dept\_id, salary)\\
departments(id, name)\\
orders(id, cust\_id, amount, created\_at)\\
order\_items(id, order\_id, product\_id, qty)\\
products(id, name, price)
}

%──────────────────────────────────────────────────────
\section{SQL задачки}

\hl{blue}{1. Средняя ЗП по отделам, только > 80k}
\begin{lstlisting}[style=sqlcode]
SELECT d.name, AVG(e.salary) avg_sal
FROM employees e
JOIN departments d ON d.id = e.dept_id
GROUP BY d.name
HAVING AVG(e.salary) > 80000
ORDER BY avg_sal DESC;
\end{lstlisting}

\hl{blue}{2. Топ-3 клиента по сумме заказов}
\begin{lstlisting}[style=sqlcode]
SELECT cust_id, SUM(amount) total
FROM orders
GROUP BY cust_id
ORDER BY total DESC LIMIT 3;
\end{lstlisting}

\hl{blue}{3. Сотрудники без отдела (LEFT JOIN)}
\begin{lstlisting}[style=sqlcode]
SELECT e.name FROM employees e
LEFT JOIN departments d ON d.id = e.dept_id
WHERE d.id IS NULL;
\end{lstlisting}

\hl{blue}{4. Выручка за 7 дней по дням}
\begin{lstlisting}[style=sqlcode]
SELECT DATE(created_at) day,
       SUM(amount) revenue
FROM orders
WHERE created_at >= NOW()-INTERVAL '7 days'
GROUP BY day ORDER BY day;
\end{lstlisting}

\hl{blue}{5. Ранг ЗП внутри отдела (window)}
\begin{lstlisting}[style=sqlcode]
SELECT name, salary,
  RANK() OVER (
    PARTITION BY dept_id
    ORDER BY salary DESC
  ) rnk
FROM employees;
\end{lstlisting}

\hl{blue}{6. Товары без заказов}
\begin{lstlisting}[style=sqlcode]
-- NOT IN опасен при NULL; лучше:
SELECT p.id, p.name FROM products p
LEFT JOIN order_items oi
  ON oi.product_id = p.id
WHERE oi.id IS NULL;
\end{lstlisting}

\vfill
\textcolor{border}{\rule{\columnwidth}{0.4pt}}\\[0.5pt]
{\fontsize{5.8}{7}\color{muted}
  \textbf{Правила:} wg.Add \emph{до} go\,·\,канал закрывает \emph{отправитель}\,·\,Mutex нереентерабелен\,·\,%
  -race обязателен\,·\,HAVING фильтрует агрегаты\,·\,оконные функции не убирают строки%
}

\end{multicols}

\end{document}
